% Fresh explanatory draft. Not integrated into the manuscript.
% Mathematical source: formal/hko-feasible-section-upper-branches.tex.
% This passage assumes capacity, sys and the HKO body were already introduced.
\subsection*{Touching upper bounds and the local maximum}

The HKO body has systolic ratio greater than one. We now ask what happens
when its ten supporting hyperplanes move. These perturbations include bodies
that are no longer products: each of the ten normal vectors can move in all
four coordinate directions. The result is local maximality among nearby
polytopes with exactly ten facets, with equality precisely along the local
family obtained by translations, positive dilations and linear symplectic
maps. These transformations preserve the ratio, so they account for the
equality cases.

Write the nearby bodies as
\[
 K(a)=\{x\in\mathbb R^4:a_i\cdot x\leq1,\quad i=1,\ldots,10\},
 \qquad a=(a_1,\ldots,a_{10}),
\]
and let $a_0$ denote the HKO normals. We restrict $a$ to a neighborhood in
which the ten inequalities remain facets and the simple face structure is
unchanged. At the end of the proof we justify why every sufficiently close
ten-facet polytope has such a representation, using polarity to label its
facets near those of HKO. The parameter space has dimension forty. Our proof constructs
smooth functions $U_1,\ldots,U_{26}$ satisfying
\[
 \operatorname{sys}(K(a))\leq U_j(a),\qquad
 U_j(a_0)=\operatorname{sys}(K(a_0)).
\]
Thus each function touches the ratio from above at the HKO body. To prove
that a nearby body has a smaller ratio, it suffices to find one of these
upper bounds that has decreased. Different bounds may handle different
directions of motion.

Here is the elementary mechanism. Suppose first that a function $f$ on
$\mathbb R^d$ has finitely many smooth upper bounds $U_j$ that agree with it
at the origin. Write $r_j=DU_j(0)$. Assume that these linear functionals span
$(\mathbb R^d)^*$ and that
\[
 \sum_j\lambda_j r_j=0,\qquad
 \lambda_j>0,\qquad \sum_j\lambda_j=1.
\]
For every nonzero direction $h$, at least one $r_j(h)$ must be negative.
Indeed, if all were nonnegative, their positively weighted sum could vanish
only if they were all zero. The spanning assumption would then force $h=0$.

This observation also gives a uniform decrease. The continuous function
$h\mapsto\min_j r_j(h)$ is negative on the unit sphere. Compactness supplies
$\eta>0$ such that its value is at most $-\eta$ everywhere on that sphere.
Since there are finitely many upper bounds, their differentiability gives a
single neighborhood in which
\[
 |U_j(v)-U_j(0)-r_j(v)|\leq\tfrac12\eta\|v\|
 \quad\text{for every }j.
\]
For nonzero $v$ in this neighborhood, choose a bound with
$r_j(v)\leq-\eta\|v\|$. Then
\[
 f(v)\leq U_j(v)\leq f(0)-\tfrac12\eta\|v\|<f(0).
\]
The uniform estimate is what turns a statement about directions into a
local maximum throughout a neighborhood.

At HKO, the derivatives cannot span all forty directions: the ratio is
constant under its continuous symmetries. Their tangent space $T$ is generated by four translation directions, one
dilation direction and ten linear symplectic directions. The exact calculation
checks that these fifteen vectors are independent; group dimension alone
would not exclude a stabilizer at this body. The finite calculation
shows that the twenty-six rows $DU_j(a_0)$ vanish on $T$, have rank
twenty-five and admit the strictly positive relation above. Choose a linear
complement $S$ to $T$. Its dimension is twenty-five, and the restricted rows
span $S^*$. The preceding argument proves a strict local maximum on the
affine slice $a_0+S$.

To recover all nearby parameters, apply a small symmetry to this slice.
The derivative of this operation at $a_0$ combines the independent spaces
$T$ and $S$, so it is invertible. The inverse function theorem therefore
expresses each nearby parameter as a symmetry of a nearby slice point.
Invariance reduces its ratio to the ratio on the slice. It is smaller unless
the slice point is $a_0$, which gives the asserted equality cases.

It remains to construct the touching upper bounds and establish the finite
rank and positivity statements. The variational capacity formula supplies
the bounds: any feasible choice of weights with positive quadratic value
$Q$ gives $c_{\mathrm{EHZ}}\leq1/(2Q)$. We choose weights that attain the
known capacity at $a_0$, then hold selected weights fixed and solve the five linear equations for
closure and normalization for the others. An invertible five-column minor
makes this solution smooth, and strict positivity of the base weights makes
positivity persist nearby. This constructs feasible weights even when the
equations for a stationary optimum are singular. The explicit seven-facet example below
shows how this continuation works and how to differentiate its upper bound.

For an ordered list $\sigma=(\sigma_1,\ldots,\sigma_m)$ of distinct facets,
write
\[
 C_\sigma(a)=
 \begin{pmatrix}a_{\sigma_1}&\cdots&a_{\sigma_m}\\1&\cdots&1\end{pmatrix},
 \qquad e=(0,0,0,0,1)^T.
\]
Its first four rows impose $\sum_i\beta_i a_{\sigma_i}=0$ and its last
row imposes $\sum_i\beta_i=1$. Thus $C_\sigma\beta=e$ and $\beta_i>0$
are precisely the closure, normalization and positivity conditions we use.
Along any smooth feasible section with positive quadratic value, put
\[
 A_\sigma(a)=\frac{1}{2Q_\sigma(a,\beta(a))},\qquad
 U_\sigma(a)=\frac{A_\sigma(a)^2}{2V(a)},\qquad
 V(a)=\operatorname{vol}_4 K(a).
\]
The variational capacity inequality and positivity give
$\operatorname{sys}(K(a))\leq U_\sigma(a)$. The exact calculations below
establish equality at $a_0$ and determine the derivatives of these functions.
